\documentclass[11pt]{amsart}

% Encoding and fonts
\usepackage[T1]{fontenc}
\usepackage[utf8]{inputenc}
\usepackage{lmodern}

% Math packages
\usepackage{amsmath}
\usepackage{amssymb}
\usepackage{amsfonts}
\usepackage{amsthm}
\usepackage{mathtools}

% Layout and utilities
\usepackage{geometry}
\usepackage{enumitem}
\usepackage{booktabs}
\usepackage{graphicx}
\geometry{margin=1in}

\numberwithin{equation}{section}
\setcounter{tocdepth}{3}
\setcounter{secnumdepth}{3}

% Theorem-like environments
\theoremstyle{plain}
\newtheorem{theorem}{Theorem}[section]
\newtheorem{proposition}[theorem]{Proposition}
\newtheorem{lemma}[theorem]{Lemma}
\newtheorem{corollary}[theorem]{Corollary}
\newtheorem{conjecture}[theorem]{Conjecture}

\theoremstyle{definition}
\newtheorem{definition}[theorem]{Definition}
\newtheorem{experiment}[theorem]{Experiment}
\newtheorem{heuristic}[theorem]{Heuristic}
\newtheorem{question}[theorem]{Question}

\theoremstyle{remark}
\newtheorem{remark}[theorem]{Remark}
\newtheorem{warning}[theorem]{Warning}

\newcommand{\mob}{\mu}
\newcommand{\Mert}{M}
\newcommand{\Sq}{Q}
\newcommand{\PS}{\mathcal P}

\title{An Ordinal-Prime Transform of the M\"obius Function}
\author{}
\date{Version 5}

% hyperref must be loaded last
\usepackage[hidelinks]{hyperref}
\begin{document}

\maketitle
\tableofcontents

\section{Purpose and scope}
\label{sec:purpose}

This note develops the ordinal-prime transform of the M\"obius function and
determines exactly how much arithmetic information it carries. It is
self-contained and uses no other construction.

The starting point is the M\"obius function $\mob(n)\in\{-1,0,+1\}$. The
occurrences of $-1$ and of $+1$ are numbered separately, in increasing order of
$n$, and the corresponding ordinal prime is attached to each occurrence; zeros
are mapped to $1$. The resulting signed sequence is $S(n)$ and its partial sum
is $SS(n)=\sum_{k\le n}S(k)$.

The conclusions are:

\begin{enumerate}[label=(\roman*)]
\item $SS(n)$ has an exact closed form in terms of the Mertens function
$\Mert(n)$, the squarefree counting function $\Sq(n)$, and the sum of the first
$m$ primes (Proposition~\ref{prop:closed}).
\item After subtracting the deterministic zero contribution $n-\Sq(n)$, the
remainder $T(n):=SS(n)-n+\Sq(n)$ is \emph{exactly} a signed sum of $|\Mert(n)|$
consecutive primes, occupying a window of prime indices symmetric about
$(\Sq(n)+1)/2$ (Proposition~\ref{prop:block}).
\item Consequently $T(n)=\Mert(n)\,p_{\lfloor \Sq(n)/2\rfloor}(1+\theta(n))$
with $\theta(n)\to0$ effectively, and the Mertens criterion transfers to $T$ in
both directions (Theorem~\ref{thm:RH}).
\item Every object built here --- $SS$, $T$, the residual $E$, and the centered
residual $E_{\mathrm{mid}}$ --- is a deterministic function of the pair
$(\Mert(n),\Sq(n))$ alone. The transform therefore creates no independent
arithmetic constraint on $\Mert(n)$. This applies to the residuals as well and
is stated as Proposition~\ref{prop:info}.
\item The nonlinearity of the transform splits exactly into a central part
linear in $\Mert(n)$ and an \emph{antisymmetric weighted sum of prime gaps}
with triangular weights (Theorem~\ref{thm:gaps}). The weight norm has a closed
form (Lemma~\ref{lem:w2}), which forces the exponent $3/2$ in the observed law
$|E_{\mathrm{mid}}(n)|\approx0.155\,|\Mert(n)|^{3/2}\log n$.
\item The measured size of that residual is quantitatively accounted for by the
local covariance structure of prime gaps. A covariance-corrected model predicts
its standard deviation to within $0.3\%$ with no free parameter, whereas an
independent-gap model underpredicts by a factor $0.853$
(Section~\ref{sec:numeric}, part D).
\item The packet structure visible in $E_{\mathrm{mid}}$ is an amplitude effect:
its moving RMS follows $|M(n)|^{3/2}\log n$. On the logarithmic coordinate
$u=\log n$, the spectrum of $M(e^u)e^{-u/2}$ displays the ordinates $\gamma_k$
of zeta zeros, as predicted by the classical explicit formula. This spectral
observation is a sanity check, not new information about the zeros; the finite
$u$-window has frequency resolution about $0.91$.
\item A separate test asks whether the prime-gap fluctuation factor is coupled to
the sign or zeta phase of $M$. No such coupling is detected at the present
scale. With effective sample about $1.7\times10^4$, the tests only exclude
linear effects larger than roughly $0.015$ and do not constitute a theorem of
independence.
\end{enumerate}

The core numerical pipeline was regenerated for $1\le n\le10^{6}$ where stated.
Additional follow-up measurements introduced in Version 5 are identified explicitly as
reported measurements when their originating script was not present in the supplied
package; see Section~\ref{sec:numeric}.

\begin{remark}[Version 5 audit scope]
\label{rem:v5-audit}
Version 5 restores section-based equation numbering and labels every displayed
formula.  It also corrects boundary and status issues in Version 5: the
multiplicative prime scale starts at $n=2$; the residuals are not claimed to be
monotone in $M$; prime gaps are not asserted pointwise to be asymptotic to
$\log p$; the cases $|M|\le2$ are separated in the gap formula; the
prime-counting integral includes the $a=0$ endpoint convention; all sums used
in that integral and in the $\kappa$ identity are explicitly finite; the
Montgomery--Soundararajan comparison is identified as a reference scale for
weighted prime counts rather than a theorem about the unweighted $\pi$ count;
and the shared-counter non-endpoint claim is stated as a functional statement
with an explicit pair of histories.
\end{remark}

\section{Notation and elementary lemmas}
\label{sec:notation}

Let $p_1=2,p_2=3,p_3=5,\ldots$ denote the primes in increasing order, and set
$p_0:=0$ by convention (used only to make the product forms below well defined
at $n=1$).

Write
\begin{equation}
\label{eq:basic-definitions}
\Mert(n):=\sum_{k\le n}\mob(k),
\qquad
\Sq(n):=\sum_{k\le n}\mob(k)^2 ,
\qquad
\PS(m):=\sum_{j=1}^{m}p_j ,\quad \PS(0):=0 .
\end{equation}

Define the cumulative occurrence counts
\begin{equation}
\label{eq:occurrence-counts}
N_-(n)=\sum_{k\le n}\frac{\mob(k)^2-\mob(k)}2,
\qquad
N_0(n)=\sum_{k\le n}\bigl(1-\mob(k)^2\bigr),
\qquad
N_+(n)=\sum_{k\le n}\frac{\mob(k)^2+\mob(k)}2 .
\end{equation}
Thus $N_-(n)$, $N_0(n)$, $N_+(n)$ count the occurrences of $\mob=-1$, $\mob=0$,
$\mob=+1$ up to $n$, and the counts are inclusive of $n$ itself.

\begin{lemma}[Counter identities]\label{lem:counters}
For every $n\ge1$,
\begin{equation}
\label{eq:counter-identities}
N_+(n)+N_-(n)=\Sq(n),
\qquad
N_+(n)-N_-(n)=\Mert(n),
\qquad
N_0(n)=n-\Sq(n),
\end{equation}
hence
\begin{equation}
\label{eq:counter-solutions}
N_+(n)=\frac{\Sq(n)+\Mert(n)}2,
\qquad
N_-(n)=\frac{\Sq(n)-\Mert(n)}2 .
\end{equation}
\end{lemma}

\begin{lemma}[Parity]\label{lem:parity}
$\Sq(n)\equiv \Mert(n)\pmod 2$ for every $n\ge1$. Consequently the indices
$(\Sq(n)\pm\Mert(n))/2$ occurring throughout are integers.
\end{lemma}

\begin{proof}
$\Sq=N_++N_-$ and $\Mert=N_+-N_-$, so $\Sq-\Mert=2N_-$.
\end{proof}

Lemma~\ref{lem:parity} is used silently in every subsequent formula and was
absent from version 1; without it the central identities are not well formed.

\section{The transform}
\label{sec:transform}

\begin{definition}
\label{def:transform}
The three rows are
\begin{equation}
\label{eq:transform-rows}
R_-(n)=-\frac{\mob(n)^2-\mob(n)}2\,p_{N_-(n)},
\qquad
R_0(n)=1-\mob(n)^2,
\qquad
R_+(n)=\frac{\mob(n)^2+\mob(n)}2\,p_{N_+(n)} ,
\end{equation}
with the convention $p_0=0$ (needed only at $n=1$, where $N_-(1)=0$ and the
coefficient vanishes). The ordinal-prime transform is
\begin{equation}
\label{eq:transform-piecewise}
S(n):=R_-(n)+R_0(n)+R_+(n)
=
\begin{cases}
-p_{N_-(n)},&\mob(n)=-1,\\[1mm]
1,&\mob(n)=0,\\[1mm]
+p_{N_+(n)},&\mob(n)=+1 .
\end{cases}
\end{equation}
\end{definition}

Along the occurrences of $\mob=-1$ the values are $-2,-3,-5,-7,-11,\ldots$; along
the occurrences of $\mob=+1$ they are $2,3,5,7,11,\ldots$, in a separate ordinal
count. Explicitly,
\begin{equation}
\label{eq:first-transform-values}
S(1),\ldots,S(12)=2,\,-2,\,-3,\,1,\,-5,\,3,\,-7,\,1,\,1,\,5,\,-11,\,1 .
\end{equation}

\begin{definition}
\label{def:ss-t}
$SS(n):=\sum_{k=1}^{n}S(k)$ and $T(n):=SS(n)-\bigl(n-\Sq(n)\bigr)$.
\end{definition}

\begin{proposition}[Closed form]\label{prop:closed}
For every $n\ge1$,
\begin{equation}
\label{eq:closed-form}
\boxed{\begin{aligned}
SS(n)&=\PS\bigl(N_+(n)\bigr)-\PS\bigl(N_-(n)\bigr)+N_0(n)\\
&=\PS\!\left(\frac{\Sq(n)+\Mert(n)}2\right)
-\PS\!\left(\frac{\Sq(n)-\Mert(n)}2\right)+n-\Sq(n).
\end{aligned}}
\end{equation}
\end{proposition}

\begin{proof}
The $k\le n$ with $\mob(k)=+1$ receive, in order, the primes
$p_1,\dots,p_{N_+(n)}$, contributing $\PS(N_+(n))$; those with $\mob(k)=-1$
contribute $-\PS(N_-(n))$; each of the $N_0(n)$ zeros contributes $1$. The
second form is Lemma~\ref{lem:counters}, the indices being integers by
Lemma~\ref{lem:parity}.
\end{proof}

\section{The signed prime block}
\label{sec:block}

\begin{proposition}[Block representation]\label{prop:block}
Write $W(n)$ for the set of prime indices strictly between $N_-(n)$ and
$N_+(n)$, inclusive of the larger endpoint; $W(n)$ has exactly $|\Mert(n)|$
elements and is symmetric about the index $(\Sq(n)+1)/2$. Then
\begin{equation}
\label{eq:block-representation}
T(n)=\operatorname{sgn}\bigl(\Mert(n)\bigr)\sum_{j\in W(n)}p_j
=
\begin{cases}
\displaystyle\sum_{j=(\Sq-\Mert)/2+1}^{(\Sq+\Mert)/2}p_j, & \Mert(n)>0,\\[5mm]
0,& \Mert(n)=0,\\[2mm]
\displaystyle-\sum_{j=(\Sq-|\Mert|)/2+1}^{(\Sq+|\Mert|)/2}p_j, & \Mert(n)<0 .
\end{cases}
\end{equation}
In particular
\begin{equation}
\label{eq:block-sign}
\operatorname{sgn}T(n)=\operatorname{sgn}\Mert(n),
\qquad
T(n)=0\iff \Mert(n)=0 .
\end{equation}
\end{proposition}

\begin{proof}
Immediate from Proposition~\ref{prop:closed}: $T=\PS(N_+)-\PS(N_-)$ is a
telescoping difference of prefix sums of the primes. The window endpoints are
$\min(N_\pm)+1$ and $\max(N_\pm)$, whose midpoint is
$(N_++N_-+1)/2=(\Sq+1)/2$.
\end{proof}

\begin{corollary}[Exact bounds]\label{cor:bounds}
\begin{equation}
\label{eq:block-bounds}
|\Mert(n)|\,p_{\frac{\Sq(n)-|\Mert(n)|}{2}+1}
\;\le\;
|T(n)|
\;\le\;
|\Mert(n)|\,p_{\frac{\Sq(n)+|\Mert(n)|}{2}} .
\end{equation}
\end{corollary}

Both inequalities are trivial (each of the $|\Mert(n)|$ primes in the block lies
between the smallest and the largest one), but the \emph{lower} bound is what
makes the transfer of Mertens bounds an equivalence rather than an implication;
see Theorem~\ref{thm:RH}.

\section{The prime scale, with explicit error}
\label{sec:prime-scale}

We record the multiplicative estimate in a form that is valid at every $n$,
including all $n\ge2$ with $\Mert(n)=0$. The case $n=1$ is excluded only because
$\lfloor\Sq(1)/2\rfloor=0$ and the reference prime $p_0=0$ cannot be used as a
multiplicative scale. (Version 1 wrote
$T(n)\sim \Mert(n)p_{\Sq(n)/2}$; that assertion is empty on the zero set of
$\Mert$, which contains $5361$ values of $n\le10^{6}$.)

\begin{proposition}\label{prop:scale}
For $n\ge2$, set $\theta(n)=0$ when $M(n)=0$. When $M(n)\ne0$, define
$\theta(n)$ by equation~\eqref{eq:theta-definition}. Then
\begin{equation}
\label{eq:theta-definition}
T(n)=M(n)\,p_{\lfloor Q(n)/2\rfloor}\bigl(1+\theta(n)\bigr).
\end{equation}
\begin{equation}
\label{eq:theta-bound}
|\theta(n)|\;\le\;
\max_{|j-\Sq(n)/2|\le |\Mert(n)|/2+1}
\left|\frac{p_j}{p_{\lfloor \Sq(n)/2\rfloor}}-1\right| ,
\end{equation}
and $\theta(n)\to0$ as $n\to\infty$.
\end{proposition}

\begin{proof}
The bound is Proposition~\ref{prop:block} together with the triangle
inequality. For the limit: by the prime number theorem in the form
$\Mert(n)=o(n)$, and by $\Sq(n)=\tfrac{6}{\pi^2}n+O(\sqrt n)$, one has
$|\Mert(n)|/\Sq(n)\to0$; hence the window endpoints are
$(1+o(1))\Sq(n)/2$ on the index scale, and $p_{(1+o(1))m}/p_m\to1$ by the prime
number theorem again.
\end{proof}

Both hypotheses invoked here are named: $\Mert(n)=o(n)$ is equivalent to the prime number theorem
\cite{Titchmarsh1986}, and $\Sq(n)=\tfrac6{\pi^2}n+O(\sqrt n)$ is the standard
squarefree-counting estimate \cite{Apostol1976}.

\begin{remark}[The asymptotic constant is a poor numerical guide]
\label{rem:prime-scale-data}
From $p_m=m(\log m+\log\log m-1+o(1))$ and $\Sq(n)/2\sim \tfrac3{\pi^2}n$ one
gets
\begin{equation}
\label{eq:prime-scale-expansion}
\begin{aligned}
p_{\lfloor \Sq(n)/2\rfloor}
&=\frac{3}{\pi^2}\,n\log n\\
&\quad\times\left(1+\frac{\log\log n-1+\log(3/\pi^2)}{\log n}
+o\!\left(\tfrac1{\log n}\right)\right),
\end{aligned}
\end{equation}
so the correction to the leading constant decays only like $1/\log n$. Measured
values are
\begin{equation}
\label{eq:leading-scale-data}
\frac{T(n)}{(3/\pi^2)n\log n\,\Mert(n)}
=0.9558,\ 0.9954,\ 1.0151,\ 1.0281
\end{equation}
at $n=10^3,10^4,10^5,10^6$; they are not monotone and remain $2.8\%$ off at
$n=10^6$. By contrast,
\begin{equation}
\label{eq:exact-prime-scale-data}
\frac{T(n)}{\Mert(n)p_{\lfloor \Sq(n)/2\rfloor}}
=1.00200,\ 1.00071,\ 1.000097,\ 0.999998
\end{equation}
at the same points. Any quantitative use should keep
$p_{\lfloor \Sq(n)/2\rfloor}$ and not the closed constant.
\end{remark}

\section{Transfer of the Mertens criterion}
\label{sec:mertens-transfer}

\begin{proposition}
\label{prop:growth}
$|T(n)|=O\bigl(n\log n\,|\Mert(n)|\bigr)$ and hence
$SS(n)=O\bigl(n\log n\,|\Mert(n)|+n\bigr)$. Unconditionally
$SS(n)=o(n^2\log n)$, by $\Mert(n)=o(n)$ (prime number theorem).
\end{proposition}

\begin{theorem}\label{thm:RH}
The following are equivalent.
\begin{enumerate}[label=(\alph*)]
\item The Riemann hypothesis.
\item $\Mert(n)=O_\varepsilon(n^{1/2+\varepsilon})$ for every $\varepsilon>0$.
\item $SS(n)-n+\Sq(n)=O_\varepsilon(n^{3/2+\varepsilon})$ for every
$\varepsilon>0$.
\end{enumerate}
\end{theorem}

\begin{proof}
(a)$\iff$(b) is the classical Littlewood criterion; see \cite{Titchmarsh1986}.

(b)$\Rightarrow$(c): fix $\varepsilon>0$ and apply (b) with
$\eta=\varepsilon/2$. By Corollary~\ref{cor:bounds} and the prime number theorem,
$|T(n)|\ll |\Mert(n)|\,n\log n\ll_{\varepsilon}
 n^{3/2+\varepsilon/2}\log n=O_{\varepsilon}(n^{3/2+\varepsilon})$.

(c)$\Rightarrow$(b): this is the direction version 1 asserted without proof. By
the prime number theorem $|\Mert(n)|=o(n)$ and $\Sq(n)=\tfrac6{\pi^2}n+O(\sqrt
n)$, so
\begin{equation}
\label{eq:smallest-block-index}
\frac{\Sq(n)-|\Mert(n)|}{2}+1=\bigl(1+o(1)\bigr)\frac{3}{\pi^2}n ,
\end{equation}
whence the smallest prime of the block satisfies
$p_{(\Sq-|\Mert|)/2+1}\ge c\,n\log n$ for some $c>0$ and all large $n$. The
lower bound of Corollary~\ref{cor:bounds} then gives
\begin{equation}
\label{eq:reverse-mertens-bound}
|\Mert(n)|\le \frac{|T(n)|}{c\,n\log n}
\ll_\varepsilon \frac{n^{3/2+\varepsilon}}{n\log n}
\ll_\varepsilon n^{1/2+\varepsilon}. \qedhere
\end{equation}
\end{proof}

\begin{remark}
\label{rem:rh-reformulation}
Theorem~\ref{thm:RH} is a reformulation, not a proof mechanism. Note also that
(b) must not be strengthened to $\Mert(n)=O(\sqrt n)$: the Riemann hypothesis
does not imply the Mertens conjecture, which is false \cite{OdlyzkoTeRiele1985}.
\end{remark}

\section{The zeta ingredients}
\label{sec:zeta-ingredients}

The two M\"obius-derived ingredients have the classical Dirichlet series \cite{Apostol1976}
\begin{equation}
\label{eq:dirichlet-series}
\sum_{m\ge1}\frac{\mob(m)}{m^s}=\frac1{\zeta(s)},
\qquad
\sum_{m\ge1}\frac{\mob(m)^2}{m^s}=\frac{\zeta(s)}{\zeta(2s)},
\qquad \Re s>1,
\end{equation}
and the prime weights are governed by $\log\zeta(s)$ through the prime-counting
function, via the exact Abel identity
\begin{equation}
\label{eq:abel-prime-sum}
\PS(m)=m\,p_m-\int_2^{p_m}\pi(t)\,dt ,
\end{equation}
(verified here for $m\le100$; e.g.\ $\PS(3)=10=3\cdot5-5$).

Version 1 substituted Perron integrals for $\Mert$ and $\Sq$ into the arguments
of $\PS$. That is not a functional identity and it is dropped here, for two
reasons. First, $\PS$ is defined only on nonnegative integers, so an integral
placed inside its argument is notation for ``the integer that this integral
evaluates to'', not an analytic representation. Second, at integer $x$ the
untruncated Perron integral returns $\Mert(x)-\tfrac12\mob(x)$, not $\Mert(x)$.

The honest statement is the following. Equations~\eqref{eq:dirichlet-series} and
\eqref{eq:abel-prime-sum} show that the transform composes three classical objects,
\begin{equation}
\label{eq:zeta-components}
\frac1{\zeta(s)}\ \ (\text{signed M\"obius imbalance}),
\qquad
\frac{\zeta(s)}{\zeta(2s)}\ \ (\text{squarefree count}),
\qquad
\log\zeta(s)\ \ (\text{prime weights}),
\end{equation}
through the nonlinear expression
$\PS\bigl(\tfrac{\Sq+\Mert}2\bigr)-\PS\bigl(\tfrac{\Sq-\Mert}2\bigr)+n-\Sq$.
None of the three components is new, and the composition is not a Dirichlet
series of anything.

\section{Exact information content}
\label{sec:information}

\begin{proposition}\label{prop:info}
For each $n$, $SS(n)$ is a deterministic function of the pair
$(\Mert(n),\Sq(n))$. With $\Sq(n)$ held fixed and $\Mert(n)$ restricted to the
admissible parity class $\Mert(n)\equiv\Sq(n)\pmod 2$, $T(n)$ is strictly
increasing in $\Mert(n)$. Every residual defined below, including $E(n)$ and
$E_{\mathrm{mid}}(n)$, is likewise a deterministic function of
$(\Mert(n),\Sq(n))$; no monotonicity claim is made for the residuals.
\end{proposition}

\begin{proof}
The first claim is Lemma~\ref{lem:counters} plus
Proposition~\ref{prop:closed}. For the second, increasing $\Mert$ by $2$ at
fixed $\Sq$ raises $N_+$ by one and lowers $N_-$ by one, changing $T$ by
$p_{N_++1}+p_{N_-}>0$. The residuals $E$ and $E_{\mathrm{mid}}$ are $T$ minus explicit functions of
$(\Mert,\Sq)$, hence are also functions of $(\Mert,\Sq)$. This last statement
is only a dependence statement; subtracting the central term does not preserve the
preceding monotonicity automatically.
\end{proof}

Two consequences deserve to be stated plainly.

First, $SS(n)$ does not retain the internal ordering of the $\pm1$ values of
$\mob$ below $n$: two M\"obius sign histories with the same endpoint counts
produce the same $SS(n)$. It is a nonlinear endpoint statistic.

Second --- and this corrects a claim in version 1 --- the residual $E(n)$ is
\emph{not} structurally distinct from $SS(n)$ in the informational sense. Its
window is determined by $(\Mert(n),\Sq(n))$, so the argument above applies to it
verbatim. What is genuinely different about the residual is not its information
content but its \emph{size}, which turns out to be governed by prime-gap
statistics; that is the subject of the next two sections.

There are two distinct meanings of ``information'' that must not be mixed.
In the logical sense, Proposition~\ref{prop:info} is final:
$E_{\mathrm{mid}}(n)=G(\Mert(n),\Sq(n))$ for an explicit deterministic $G$, so
no experiment can make $E_{\mathrm{mid}}$ contain arithmetic information absent
from the endpoint pair. In the analytic or \emph{leverage} sense, one can still
ask whether a separately proved bound for the prime-gap fluctuation factor could
combine with the identity to constrain $\Mert$. A bound for $T$ itself cannot do
this: at fixed $\Sq$, $T$ is strictly monotone in $\Mert$ and the two-sided
block bounds make a Mertens bound and the corresponding $T$ bound equivalent up
to the explicit prime scale. What remains testable is only whether some
independently controlled feature of the gap fluctuation is coupled to the sign,
phase, or trajectory of $\Mert$. Section~\ref{sec:coupling} measures a small
family of such couplings and finds none at the present resolution; this is a
negative numerical result, not a logical strengthening of
Proposition~\ref{prop:info}.

\section{The residual and its correct centering}
\label{sec:centering}

\begin{definition}
\label{def:residuals}
Let
\begin{equation}
\label{eq:pstar-definition}
p^{*}(n):=\tfrac12\Bigl(p_{\lfloor (\Sq(n)+1)/2\rfloor}+p_{\lceil (\Sq(n)+1)/2\rceil}\Bigr)
\end{equation}
be the linear interpolation of the prime sequence at the exact window centre
$(\Sq(n)+1)/2$ of Proposition~\ref{prop:block}. Define
\begin{equation}
\label{eq:residual-definitions}
E(n):=T(n)-\Mert(n)\,p_{\lfloor \Sq(n)/2\rfloor},
\qquad
E_{\mathrm{mid}}(n):=T(n)-\Mert(n)\,p^{*}(n),
\end{equation}
so that the exact decomposition
\begin{equation}
\label{eq:centered-decomposition}
\boxed{\;T(n)=\Mert(n)\,p^{*}(n)+E_{\mathrm{mid}}(n)\;}
\end{equation}
holds for every $n$.
\end{definition}

When $\Sq(n)$ is odd, $(\Sq+1)/2$ is an integer and $p^{*}=p_{(\Sq+1)/2}$; when
$\Sq(n)$ is even, $p^{*}$ is the mean of the two central primes.

\begin{proposition}[Centering offset]\label{prop:offset}
$E(n)-E_{\mathrm{mid}}(n)=\Mert(n)\bigl(p^{*}(n)-p_{\lfloor \Sq(n)/2\rfloor}\bigr)$,
where, for $\Sq(n)\ge2$, the bracket equals one half of a prime gap when
$\Sq(n)$ is even and one full prime gap when $\Sq(n)$ is odd. At $n=1$ the
identity still holds under $p_0=0$, with bracket $p_1-p_0=2$. It is positive. Pointwise it is exactly one half or one full prime gap; no
uniform two-sided estimate $\asymp\log p^{*}(n)$ is asserted. Its empirical
typical scale in the computed range is comparable with $\log p^{*}(n)$, so $E$
carries a deterministic component linear in $\Mert(n)$ which
$E_{\mathrm{mid}}$ does not.
\end{proposition}

The identity was verified exactly for all $n\le10^{6}$; the median of
$\bigl(p^{*}-p_{\lfloor \Sq/2\rfloor}\bigr)/\log p_{\lfloor \Sq/2\rfloor}$ over
$10^{4}\le n\le10^{6}$ is $0.5318$. Proposition~\ref{prop:offset} explains the
correlation structure of the two residuals: the offset is a fixed positive
multiple of $\Mert(n)$, so $E$ correlates with $\Mert$ and $E_{\mathrm{mid}}$
does not. Measured over $100\le n\le10^{6}$,
\begin{equation}
\label{eq:residual-correlations}
\operatorname{corr}(E,\Mert)=+0.2573,
\qquad
\operatorname{corr}(E_{\mathrm{mid}},\Mert)=+0.0045 .
\end{equation}
$E_{\mathrm{mid}}$ is therefore the object to study; $E$ is $E_{\mathrm{mid}}$
plus a bookkeeping artefact. Version 1 studied only $E$.

\section{Exact antisymmetric gap representation}
\label{sec:gap-representation}

Write $g_i:=p_{i+1}-p_i$ and $m(n):=\lfloor(\Sq(n)+1)/2\rfloor$.

\begin{theorem}[Gap representation]\label{thm:gaps}
Let $h:=|\Mert(n)|$. If $h\le2$, then $E_{\mathrm{mid}}(n)=0$. For $h\ge3$,
\begin{equation}
\label{eq:gap-representation}
\boxed{\;
E_{\mathrm{mid}}(n)=\operatorname{sgn}\bigl(\Mert(n)\bigr)\sum_{i}w_i\,g_i \;}
\end{equation}
where the weights $w_i$ are supported on finitely many gaps symmetrically
placed about $m=m(n)$, are antisymmetric in sign about the centre, and are given
explicitly by
\begin{equation}
\label{eq:even-weights}
h \text{ even}:\quad
w_{m+k}=\tfrac h2-k,\quad w_{m-k}=-\bigl(\tfrac h2-k\bigr),\quad 1\le k\le \tfrac h2-1;
\end{equation}
\begin{equation}
\label{eq:odd-weights}
h \text{ odd}:\quad
w_{m+k}=\tfrac{h-1}2-k \ \ (0\le k\le \tfrac{h-3}2),
\qquad
w_{m-k}=-\bigl(\tfrac{h+1}2-k\bigr)\ \ (1\le k\le \tfrac{h-1}2).
\end{equation}
For even $h$, the single gap $g_m$ between the two central primes cancels. For odd
$h$, the centre is the prime $p_m$ and there is no single central gap; the two adjacent
gaps belong to the two antisymmetric sides. In either parity case the nonzero weight
magnitudes on each side are $\{1,2,\dots,\lfloor (h-1)/2\rfloor\}$. The cases
$h=0,1,2$ are the zero-residual cases stated above.
\end{theorem}

\begin{proof}
The cases $h=0,1,2$ follow directly from the definition: the block is empty, a
single central prime, or the two central primes whose mean is $p^*$. Hence
$E_{\mathrm{mid}}=0$. Assume now $h\ge3$ and take $\Mert(n)>0$; the case
$\Mert(n)<0$ follows by a global sign change.

Suppose $h$ is even, so $\Sq$ is even by Lemma~\ref{lem:parity}, $m=\Sq/2$ and
$p^{*}=\tfrac12(p_m+p_{m+1})$. The window is $\{m-\tfrac h2+1,\dots,m+\tfrac h2\}$.
For $j>m$, $p_j-p^{*}=\tfrac12 g_m+\sum_{i=m+1}^{j-1}g_i$; for $j\le m$,
$p_j-p^{*}=-\tfrac12 g_m-\sum_{i=j}^{m-1}g_i$. Summing over the window, the
$\tfrac12 g_m$ terms occur $\tfrac h2$ times with each sign and cancel. Counting
how many $j$ contain a given $g_{m+k}$ gives the coefficient $\tfrac h2-k$, and
symmetrically $-(\tfrac h2-k)$ for $g_{m-k}$.

If $h$ is odd then $\Sq$ is odd, $m=(\Sq+1)/2$ and $p^{*}=p_m$, the window being
$\{m-\tfrac{h-1}2,\dots,m+\tfrac{h-1}2\}$. The same count gives the stated
coefficients, the two sides again carrying magnitudes
$\tfrac{h-1}2,\dots,1$.
\end{proof}

\begin{lemma}[Weight norm]\label{lem:w2}
\begin{equation}
\label{eq:weight-norm}
\boxed{\;
\sum_i w_i^{2}=
\begin{cases}
\dfrac{h(h^{2}-1)}{12}, & h \text{ odd},\\[3mm]
\dfrac{h(h-1)(h-2)}{12}, & h \text{ even},
\end{cases}
\qquad\text{so}\qquad
\sum_i w_i^{2}=\frac{h^{3}}{12}\bigl(1+O(h^{-1})\bigr). }
\end{equation}
\end{lemma}

\begin{proof}
By Theorem~\ref{thm:gaps} each side contributes
$\sum_{j=1}^{\lfloor (h-1)/2\rfloor}j^{2}$. For $h$ odd this is
$2\cdot\frac{(\frac{h-1}2)(\frac{h+1}2)h}{6}=\frac{h(h^{2}-1)}{12}$; for $h$ even
it is $2\cdot\frac{(\frac h2-1)(\frac h2)(h-1)}{6}=\frac{h(h-1)(h-2)}{12}$.
\end{proof}

Theorem~\ref{thm:gaps}, equation~\eqref{eq:gap-representation}, and Lemma~\ref{lem:w2} are exact and unconditional. They
are the reason the exponent $3/2$ appears below: it is forced by the triangular
weights, not fitted.

\begin{remark}
\label{rem:gap-information}
Theorem~\ref{thm:gaps} does \emph{not} evade Proposition~\ref{prop:info}: the
window position $m(n)$ and its length $h$ are determined by
$(\Mert(n),\Sq(n))$, and the primes are deterministic, so $E_{\mathrm{mid}}$
remains a function of $(\Mert(n),\Sq(n))$. What the theorem changes is what the
\emph{size} of $E_{\mathrm{mid}}$ is controlled by, namely the local second-order
statistics of prime gaps.
\end{remark}

\begin{proposition}[Prime-counting integral form]\label{prop:pi-integral}
Let $h=|\Mert(n)|$, $a=\min(N_-(n),N_+(n))$ and
$b=\max(N_-(n),N_+(n))=a+h$. Retain the convention $p_0=0$ and extend
$\pi(t)$ by $\pi(t)=0$ for $0\le t<2$. Then, for $h>0$,
\begin{equation}
\label{eq:prime-counting-integral}
\boxed{
E_{\mathrm{mid}}(n)=\operatorname{sgn}(\Mert(n))
\left\{\int_{p_a}^{p_b}\bigl(b-\pi(t)\bigr)\,dt
-h\bigl(p^*(n)-p_a\bigr)\right\}. }
\end{equation}
For $h=0$ both sides are $0$ by convention.
\end{proposition}

\begin{proof}
For $j=0,\ldots,b-1$, on $[p_j,p_{j+1})$ one has $\pi(t)=j$ (with the
preceding convention when $j=0$). The sum below is finite, so there is no
interchange of infinite sums or limiting operations:
\begin{equation}
\label{eq:integral-summation}
\int_{p_a}^{p_b}(b-\pi(t))\,dt
 =\sum_{j=a}^{b-1}(b-j)(p_{j+1}-p_j)
 =\sum_{j=a+1}^{b}(p_j-p_a).
\end{equation}
The last expression is the unsigned prime block minus $h p_a$.
Subtracting $h(p^*-p_a)$ gives the unsigned block minus $hp^*$, and the global
factor $\operatorname{sgn}(\Mert)$ gives $E_{\mathrm{mid}}$.
\end{proof}

Thus each value of $E_{\mathrm{mid}}$ may also be read as an integrated prime
counting discrepancy over the finite interval $[p_a,p_b]$, of exact length
$p_b-p_a$. In the numerical range this length is typically on the scale
$h\log p_a$; no uniform pointwise asymptotic for such short blocks is asserted. This is genuine information about prime gaps in that specific
window, even though the choice of the window itself is already fixed by
$(\Mert,\Sq)$ and hence cannot create a new restriction on $\Mert$.

\section{Two variance models}
\label{sec:variance-models}

Theorem~\ref{thm:gaps} suggests modelling $E_{\mathrm{mid}}$ as a weighted sum
of gap fluctuations. Two models are compared; both are heuristics, and both were
tested against the data of Section~\ref{sec:numeric}.

\begin{definition}[Independent-gap model]
\label{def:independent-gap}
For $h\ge3$, treat the finitely many $g_i$ carrying nonzero weights near index
$m(n)$ as uncorrelated with common local variance
$\sigma_g(n)^{2}$. Then by Lemma~\ref{lem:w2}
\begin{equation}
\label{eq:sd-independent}
\operatorname{SD}_{\mathrm{ind}}(n):=\sigma_g(n)\sqrt{\textstyle\sum_i w_i^{2}}
\;\approx\;\frac{\sigma_g(n)}{\sqrt{12}}\,h^{3/2}.
\end{equation}
\end{definition}

\begin{definition}[Covariance-corrected model]
\label{def:covariance-gap}
For $h\ge3$, let $\rho(d)$ be the local autocorrelation used in the model,
normalised by $\rho(0)=1$, and put
$R_{ab}=\rho(|i_a-i_b|)$ on the support of $w$. Define
\begin{equation}
\label{eq:covariance-model}
\kappa(h):=\frac{w^{\!\top}Rw}{w^{\!\top}w},
\qquad
\operatorname{SD}_{\mathrm{cov}}(n):=\sigma_g(n)\sqrt{\kappa(h)\textstyle\sum_i w_i^{2}} .
\end{equation}
\end{definition}

For fixed $h$ the correction factor contains no statistic beyond the gap
autocorrelation itself. If
\begin{equation}
\label{eq:cd-definition}
c_d(h):=\sum_{|i-j|=d}w_iw_j,\qquad c_0(h)=\sum_iw_i^2,
\end{equation}
then exactly
\begin{equation}\label{eq:kappa-linear}
\boxed{\qquad
\kappa(h)=\frac{1}{c_0(h)}\sum_{d=0}^{D(h)}c_d(h)\rho(d),
\qquad}
\end{equation}
where $D(h)$ is the largest lag between two nonzero weights. The sum is finite.
Thus $\kappa$ is the autocorrelation sequence viewed through a family of
triangular antisymmetric kernels. Numerically,
the regenerated exact-extension script gives
$\max|\kappa_{\rm direct}-(C\rho)/c_0|<2.0\times10^{-15}$ in double
precision. Solving the finite
linear system in the reverse direction reproduces the first gap correlations
\begin{equation}
\label{eq:rho-recovered-values}
-0.0463,\ -0.0225,\ -0.0132,\ -0.0057,\ -0.0069,\ -0.0061,\ -0.0060,\ -0.0064.
\end{equation}
The maximum error reported for $d\le60$ is $6.8\times10^{-4}$. The inversion is severely
ill-conditioned (condition number about $10^8$), so in practice only the first
lags are stably recoverable. This reinforces the information statement:
$\kappa(h)$ is a compressed linear image of $\rho(d)$, not an additional
spectral observable.

Because $w$ is antisymmetric about the centre, a \emph{positive} correlation
between gaps on opposite sides of the centre enters $w^{\!\top}Rw$ with a
negative sign. One therefore expects $\kappa(h)<1$, i.e.\ the independent model
to \emph{overpredict}. This is what is observed, and it accounts quantitatively
for the discrepancy (Section~\ref{sec:numeric}, part D).

\section{Numerical verification and measured laws}
\label{sec:numeric}

\subsection{Protocol}
\label{sec:num-protocol}
$\mob$ was computed by a linear sieve for $n\le N=10^{6}$; $M,Q$ by prefix sums;
a prime table to index $304\,071$ ($p=4\,318\,813$); $T$ by prefix sums of the
primes. At $N=10^{6}$: $\Mert(N)=212$, $\Sq(N)=607\,926$, maximum required prime
index $304\,069$. All scripts and raw outputs accompany this note.

\subsection{A. Exact statements verified}
\label{sec:num-exact}
The following hold with no exceptions for all $1\le n\le10^{6}$:
Lemma~\ref{lem:parity}; Proposition~\ref{prop:closed} (direct summation of $S$
against the closed form); Proposition~\ref{prop:block} and both inequalities of
Corollary~\ref{cor:bounds}; $\operatorname{sgn}T=\operatorname{sgn}\Mert$ and
$T=0\iff \Mert=0$ (there are $5\,361$ values of $n\le10^{6}$ with $\Mert(n)=0$);
Proposition~\ref{prop:offset}. Theorem~\ref{thm:gaps} was checked on $4\,000$
random $n\in[10^{3},10^{6}]$ with maximum absolute error $0$, and
Lemma~\ref{lem:w2} for all $h\le60$ and $h\in\{100,101,300,301,370\}$.

\subsection{B. Global residual scale, \texorpdfstring{$100\le n\le10^{6}$}{100 to one million}}
\label{sec:num-global}
\begin{equation}
\label{eq:global-residual-table}
\begin{array}{lcc}
\toprule
& E & E_{\mathrm{mid}}\\
\midrule
\text{median }|\cdot| & 1.09\times10^{3} & 9.74\times10^{2}\\
\text{r.m.s.} & 5.295\times10^{3} & 5.066\times10^{3}\\
\max|\cdot| & 6.2766\times10^{4} & 5.7654\times10^{4}\\
\text{median }|\cdot|/|T| & 1.069\times10^{-5} & 8.894\times10^{-6}\\
\bottomrule
\end{array}
\end{equation}

\begin{figure}[htbp]
\centering
\includegraphics[width=.85\textwidth]{figures/fig1_Emid.png}
\caption[Centered residual up to one million]{$E_{\mathrm{mid}}(n)$ for $n\le10^{6}$.}
\label{fig:residual}
\end{figure}

\subsection{C. The \texorpdfstring{$h^{3/2}$}{h to the three halves} law and its stationarity}
\label{sec:num-scaling}
Put $Z_{3/2}(n):=E_{\mathrm{mid}}(n)/\bigl(|\Mert(n)|^{3/2}\log n\bigr)$.
Conditioning on $h=|\Mert(n)|$ over $n\ge10^{3}$:
\begin{equation}
\label{eq:h-binned-scaling}
\begin{array}{lcccccccc}
\toprule
h\in & [8,12) & [12,19) & [19,31) & [31,51) & [51,84) & [84,138) & [138,225) & [225,400)\\
\midrule
\operatorname{median}|Z_{3/2}| & 0.1426 & 0.1477 & 0.1589 & 0.1589 & 0.1556 & 0.1546 & 0.1507 & 0.1550\\
\text{count} & 36\,451 & 58\,125 & 86\,041 & 127\,696 & 152\,059 & 164\,880 & 231\,986 & 65\,791\\
\bottomrule
\end{array}
\end{equation}
and conditioning on $n$ instead:
\begin{equation}
\label{eq:n-binned-scaling}
\begin{array}{lccc}
\toprule
n\in & \operatorname{RMS}Z_{3/2} & \operatorname{median}|Z_{3/2}| & \text{count}\\
\midrule
{}[1.00,1.58)\times10^{5} & 0.2331 & 0.1549 & 51\,624\\
{}[1.58,2.51)\times10^{5} & 0.2379 & 0.1574 & 82\,660\\
{}[2.51,3.98)\times10^{5} & 0.2340 & 0.1541 & 134\,286\\
{}[3.98,6.31)\times10^{5} & 0.2314 & 0.1530 & 215\,541\\
{}[6.31,10.0)\times10^{5} & 0.2324 & 0.1552 & 359\,873\\
\bottomrule
\end{array}
\end{equation}
Both tables are flat: the normalisation $|\Mert|^{3/2}\log n$ is stable across a
factor of $46$ in $h$ and a factor of $10$ in $n$. This supports, but does not
prove, the hypothesis that $Z_{3/2}$ converges in distribution.

\begin{figure}[htbp]
\centering
\includegraphics[width=.48\textwidth]{figures/fig2_scaling.png}
\hfill
\includegraphics[width=.48\textwidth]{figures/fig3_stationarity.png}
\caption[Residual scaling and stationarity]{Left: $|E_{\mathrm{mid}}|$ against $|\Mert|$ with the envelope
$0.155\,h^{3/2}\log n$. Right: stationarity of $Z_{3/2}$ in $n$.}
\label{fig:scaling-stationarity}
\end{figure}

\subsection{D. Which variance model is right}
\label{sec:num-variance}
Local gap statistics were measured on the prime-index band
$[5\times10^{4},3.05\times10^{5})$, which covers the windows for
$n\in[10^{5},10^{6}]$: mean gap $14.5927$, $\log p=14.6984$ at the band centre,
$\operatorname{Var}(g)=145.93$, hence
\begin{equation}
\label{eq:gap-variance-ratio}
\operatorname{Var}(g)/\log^{2}p=0.6755
\end{equation}
against the value $1$ predicted by a Poisson (Cram\'er) gap model. The gap
autocorrelation is negative at short lag and positive over a long tail:
\begin{equation}
\label{eq:gap-autocorr-table}
\begin{array}{lcccccc}
\toprule
d & 1 & 2 & 3 & 4 & 5 & 6\\
\midrule
\rho(d) & -0.0463 & -0.0225 & -0.0132 & -0.0057 & -0.0069 & -0.0061\\
\bottomrule
\end{array}
\end{equation}
with $\sum_{1\le d\le420}\rho(d)=+0.4646$. The resulting correction factor
decreases steadily,
\begin{equation}
\label{eq:kappa-table}
\begin{array}{lccccccc}
\toprule
h & 5 & 10 & 20 & 40 & 80 & 160 & 320\\
\midrule
\kappa(h) & 1.0206 & 0.9532 & 0.9094 & 0.8514 & 0.7869 & 0.7381 & 0.6871\\
\bottomrule
\end{array}
\end{equation}

Testing both models on the $883\,123$ values of $n$ with $n\ge10^{4}$ and
$h\ge12$:
\begin{equation}
\label{eq:model-comparison-table}
\begin{array}{lcccc}
\toprule
\text{model} & \operatorname{RMS}(E_{\mathrm{mid}})/\operatorname{RMS}(\mathrm{SD}) & \operatorname{RMS}Z & \operatorname{mean}Z & \operatorname{median}|Z| \\
\midrule
\text{independent gaps} & 0.8527 & 0.8963 & -1.88\times10^{-4} & 0.5962\\
\text{covariance-corrected} & 1.0031 & 1.0036 & -2.13\times10^{-4} & 0.6694\\
\bottomrule
\end{array}
\end{equation}
The $15\%$ deficit of the independent model is therefore explained, to within
$0.3\%$, by the correlations between prime gaps; and the corrected model
reproduces the Gaussian value $\operatorname{median}|Z|=0.6745$ to within $0.8\%$
with no free parameter. Within the measured range, this accounts for the scale discrepancy left by the independent model.

\begin{figure}[htbp]
\centering
\includegraphics[width=.48\textwidth]{figures/fig4_kappa.png}
\hfill
\includegraphics[width=.48\textwidth]{figures/fig5_qq.png}
\caption{Left: the correction factor $\kappa(h)$. Right: quantiles of
$Z=E_{\mathrm{mid}}/\mathrm{SD}$ against the normal, on a thinned sample; the
independent-gap curve is visibly compressed, the corrected one is not.}
\label{fig:kappa-qq}
\end{figure}

\subsection{E. Distribution of the standardised residual}
\label{sec:num-distribution}
With $Z:=E_{\mathrm{mid}}/\operatorname{SD}_{\mathrm{cov}}$, by blocks in $n$
($h\ge12$):
\begin{equation}
\label{eq:z-block-table}
\begin{array}{lcccccc}
\toprule
n\in & \text{count} & \operatorname{mean} & \operatorname{sd} & \text{skew} & \text{exc.\ kurt.} & \operatorname{median}|Z|\\
\midrule
{}[10^{4},10^{5}) & 67\,120 & +0.0043 & 0.9813 & -0.0098 & +0.2556 & 0.6474\\
{}[10^{5},3.16\times10^{5}) & 183\,261 & -0.0000 & 1.0029 & -0.0167 & +0.1559 & 0.6705\\
{}[3.16\times10^{5},10^{6}) & 632\,741 & -0.0008 & 1.0061 & -0.0188 & +0.0957 & 0.6715\\
\bottomrule
\end{array}
\end{equation}
and by blocks in $h$, the excess kurtosis decreases overall, with a small
nonmonotone rise between the fourth and fifth bins,
\begin{equation}
\label{eq:kurtosis-by-h}
\begin{array}{lccccccc}
\toprule
h\in & [12,19) & [19,31) & [31,51) & [51,84) & [84,138) & [138,225) & [225,400)\\
\midrule
\text{exc.\ kurt.} & +0.676 & +0.338 & +0.213 & +0.031 & +0.058 & +0.024 & -0.143\\
\bottomrule
\end{array}
\end{equation}
which is qualitatively consistent with, but does not establish, the behaviour expected if a central limit theorem governs the sum of
Theorem~\ref{thm:gaps} as $h\to\infty$. Quantiles of $Z$ over all $n\ge10^{4}$,
$h\ge12$:
\begin{equation}
\label{eq:quantiles-table}
\begin{array}{lccccccc}
\toprule
\text{level} & 1\% & 5\% & 25\% & 50\% & 75\% & 95\% & 99\%\\
\midrule
\text{observed} & -2.372 & -1.652 & -0.666 & +0.000 & +0.673 & +1.641 & +2.346\\
N(0,1) & -2.326 & -1.645 & -0.674 & 0 & +0.674 & +1.645 & +2.326\\
\bottomrule
\end{array}
\end{equation}

\subsection{F. The points are not independent}
\label{sec:num-dependence}
This is a necessary caveat on the preceding tables. Consecutive $n$ share almost
the same window, so the $883\,123$ values of $Z$ are strongly dependent: the
autocorrelation of $Z$ in $n$ is $+0.979$ at lag $1$, $+0.692$ at lag $10$, and
first falls below $0.05$ at lag $128$. The integrated autocorrelation time is
$\approx51$, giving an effective sample size of about $1.7\times10^{4}$, not
$8.8\times10^{5}$.

Consequently a Kolmogorov--Smirnov test on the full sample is meaningless. On
thinned subsamples:
\begin{equation}
\label{eq:thinning-table}
\begin{array}{lccccc}
\toprule
\text{thinning} & \text{count} & \operatorname{sd} & \text{exc.\ kurt.} & \mathrm{KS} & p\\
\midrule
1 & 883\,123 & 1.0036 & +0.120 & 0.00317 & 3.8\times10^{-8}\\
51 & 17\,317 & 1.0036 & +0.170 & 0.00451 & 0.87\\
128 & 6\,900 & 1.0057 & +0.018 & 0.00775 & 0.80\\
\bottomrule
\end{array}
\end{equation}
On the two thinned samples the reported KS tests do not reject a standard normal
reference law. This is descriptive rather than a proof of normality: thinning reduces
dependence but does not establish independence. The full-sample $p$-value is not
interpreted because the observations are strongly dependent.

\begin{figure}[htbp]
\centering
\includegraphics[width=.55\textwidth]{figures/fig6_autocorr.png}
\caption{Autocorrelation of $Z$ in $n$; the effective sample is roughly $50$
times smaller than the nominal one.}
\label{fig:autocorr}
\end{figure}

\subsection{G. Relative size of the residual}
\label{sec:num-relative}
Combining $|E_{\mathrm{mid}}|\approx c\,h^{3/2}\log n$ with $|T|\approx
h\,p^{*}\approx \tfrac{3}{\pi^{2}}h\,n\log n$ predicts
\begin{equation}
\label{eq:relative-residual-prediction}
\begin{aligned}
\frac{|E_{\mathrm{mid}}(n)|}{|T(n)|}
&\approx \frac{c\pi^{2}}{3}\cdot\frac{\sqrt{|\Mert(n)|}}{n},\\
\frac{c\pi^{2}}{3}&=0.510 \qquad (c=0.155).
\end{aligned}
\end{equation}
Measured medians of $\bigl(|E_{\mathrm{mid}}|/|T|\bigr)\,n/\sqrt{|\Mert|}$:
$0.4385,\ 0.4800,\ 0.4995,\ 0.4955$ on
\begin{equation}
\label{eq:relative-residual-ranges}
[10^{3},10^{4}),\quad [10^{4},10^{5}),\quad
[10^{5},3.16\times10^{5}),\quad [3.16\times10^{5},10^{6}).
\end{equation}
This quantitatively explains why the central term reproduces $T$ so accurately.
The measured median satisfies
\begin{equation}
\label{eq:median-relative-residual}
\operatorname{median}\frac{|E_{\mathrm{mid}}|}{|T|}=8.9\times10^{-6}.
\end{equation}
If $|\Mert(n)|$ is of size roughly $n^{1/2}$, the relative correction is of
size $n^{-3/4}$.

\subsection{H. Largest excursions}
\label{sec:num-extremes}
The extremes of $E_{\mathrm{mid}}$ on $n\le10^{6}$ occur in tight clusters,
e.g.\ $-5.7654\times10^{4}$ at $n=924\,067$ ($\Mert=-301$) and
$+4.6576\times10^{4}$ at $n=932\,737$ ($\Mert=-269$). These are the $n$ at which
$|\Mert|$ is locally largest, consistent with the $h^{3/2}$ law and carrying no
separate endpoint information.

\subsection{I. The packet envelope}
\label{sec:num-packets}
The visible packets are almost entirely the excursions of $|\Mert(n)|$ passed
through the $h^{3/2}$ variance scale. With a moving RMS window of width $2000$,
the log--log correlation between the measured envelope of $E_{\mathrm{mid}}$
and a moving average of $|\Mert(n)|^{3/2}\log n$ is
\begin{equation}
\label{eq:packet-correlation}
0.9939\qquad(n\ge5\times10^4).
\end{equation}
Replacing the Mertens factor by a trend depending only on $n$ gives correlation
$0.4442$. Figure~\ref{fig:packet-envelope} displays the resulting amplitude
modulation.

\begin{figure}[htbp]
\centering
\includegraphics[width=.88\textwidth]{figures/fig7_envelope.png}
\caption{The packet structure of $E_{\mathrm{mid}}$ follows the excursions of
$|M(n)|$. The red curves show a scaled moving average proportional to
$|M(n)|^{3/2}\log n$.}\label{fig:packet-envelope}
\end{figure}

\subsection{J. Spectrum on the logarithmic coordinate: a classical check}
\label{sec:num-spectrum}
Put $u=\log x$ and, for integer samples $x=n$, consider
\begin{equation}
\label{eq:log-coordinate-signal}
X(u):=e^{-u/2}\Mert(\lfloor e^u\rfloor).
\end{equation}
The classical Perron/explicit-formula treatment of the Mertens function \cite{Titchmarsh1986}, after
truncation and contour displacement and under the usual hypotheses on the zeros, has
residue contributions of the form
\begin{equation}
\label{eq:explicit-formula-mode}
\frac{x^{\rho}}{\rho\zeta'(\rho)}.
\end{equation}
For a simple zero $\rho=\beta+i\gamma$, the corresponding contribution after
$x=e^u$ and multiplication by $e^{-u/2}$ is proportional to
$e^{(\beta-1/2)u}e^{i\gamma u}$. Under RH this becomes a pure mode
$e^{i\gamma u}$. Consequently, within that classical explicit-formula framework,
the appearance of power at
frequencies $\gamma_k$ in $X(u)$ is the expected explicit-formula behaviour,
not a new law of the transform. Historical numerical uses of zeta zeros to reconstruct
prime-counting structure include \cite{RieselGohl1970}.

The numerical spectrum over $10^3\le n\le10^6$ (after the interpolation and
windowing used by the supplied script) shows peaks near the first zeta
ordinates, while a random-sign surrogate with the same squarefree support does
not show the same visual alignment. The original peak-to-nearest-zero summary
(mean distances $0.263$ for the true sequence and $1.738$ for one surrogate)
has a selection bias because peaks were selected before matching them to the
zeros. It is retained only as a diagnostic and not used as evidence. The proper
fixed-frequency question is the power evaluated at preassigned $\gamma_k$, or a
regression on the preassigned modes.

More fundamentally, the finite logarithmic window itself imposes the resolution
\begin{equation}
\label{eq:frequency-resolution}
\Delta\omega\simeq\frac{2\pi}{u_{\max}-u_{\min}}
 =\frac{2\pi}{\log(10^6/10^3)}\simeq0.91.
\end{equation}
Peak offsets of a few tenths are therefore below the resolving power of this finite
window and carry no precision information about the known $\gamma_k$. Doubling the
frequency resolution requires doubling the range in $u$, hence squaring the
range in $n$; reaching resolution $10^{-3}$ would require a logarithmic span of
order $2\pi\times10^3$, far outside this computation.

\begin{figure}[htbp]
\centering
\includegraphics[width=.88\textwidth]{figures/fig8_spectrum.png}
\caption{Power spectrum of $M(e^u)e^{-u/2}$ on $u=\log n$. Dashed lines are the
first zeta ordinates. This is an illustration of the classical explicit-formula
modes; the peak-matching statistic is not used inferentially because of
selection bias and finite-window resolution.}\label{fig:spectrum}
\end{figure}

\subsection{K. Prime-counting variance comparison}
\label{sec:num-prime-variance}
Proposition~\ref{prop:pi-integral} puts the residual on the prime-counting side.
At prime scale $X\simeq4.3\times10^6$, the supplied experiment compares the
measured variance of short-interval prime counts with an $H\log(X/H)$ reference
scale motivated by Montgomery and Soundararajan for weighted prime counts
$\psi(x+H)-\psi(x)$, and with a Poisson/Cram\'er-type $H\log X$ reference.
This is only a scale comparison. It is not a direct application of the
Montgomery--Soundararajan statement \cite{MontgomerySoundararajan2004} to the
unweighted count $\pi(x+H)-\pi(x)$. The Poisson comparison is in the spirit of
Cram\'er's model \cite{Cramer1936}. The reported ratios are
\begin{equation}
\label{eq:short-interval-variance-table}
\begin{array}{lrrrr}
\toprule
H & 200 & 1000 & 5000 & 20000\\
\midrule
\operatorname{Var}_{\rm meas}/\operatorname{Var}_{\rm MS}
 &1.116&1.202&1.310&1.622\\
\operatorname{Var}_{\rm meas}/\operatorname{Var}_{\rm Cramer}
 &0.552&0.468&0.372&0.314\\
\bottomrule
\end{array}
\end{equation}
The decreasing ratio to the Cram\'er scale is the expected qualitative signature
of an $H\log(X/H)$ variance rather than $H\log X$. The excess over the leading
MS term is not interpreted as a new statistic: at $H=2\times10^4$ the displayed
leading logarithmic factor is only about $2.96$, versus $\log X\simeq15.3$, so
lower-order terms are not negligible in this preasymptotic range.

\subsection{L. Cross-coupling test}
\label{sec:coupling}
The only numerical loophole left inside $E_{\mathrm{mid}}$ is weaker than new
endpoint information. The gap factor in Theorem~\ref{thm:gaps} is controlled by
prime statistics, while the window length is $h=|M(n)|$. One can ask whether
the standardised fluctuation
\begin{equation}
\label{eq:standardized-residual}
Z(n)=\frac{E_{\mathrm{mid}}(n)}{\operatorname{SD}_{\mathrm{cov}}(n)},
\qquad h=|M(n)|\ge12,
\end{equation}
is coupled to the sign or to a zeta phase of $M$. The supplied follow-up tests
reported
\begin{equation}
\label{eq:coupling-table}
\begin{array}{lc}
\toprule
\text{test} & \text{value}\\
\midrule
\operatorname{corr}(Z,\operatorname{sgn}M) & +0.00058\\
\operatorname{corr}(Z,M/\sqrt n) & +0.00110\\
\operatorname{corr}(Z,\cos(\gamma_1\log n)) & -0.00119\\
\operatorname{corr}(Z,\sin(\gamma_1\log n)) & +0.00173\\
\overline Z\,|_{M>0}-\overline Z\,|_{M<0} & +0.00027\quad(0.08\sigma\ \text{effective})\\
\bottomrule
\end{array}
\end{equation}
and an eight-bin phase test gave dispersion of the bin means $0.00070$, compared
with a typical standard error $0.00499$.

These are null results only at the tested resolution. With effective sample
$n_{\rm eff}\simeq17\,316$, the nominal standard error of a small correlation is
$1/\sqrt{n_{\rm eff}}\simeq0.0076$, so the experiment excludes only linear
couplings of order $|r|\gtrsim0.015$. Weaker effects, nonlinear dependence,
untested functions of the phase, or effects emerging only at much larger $n$
remain outside the sensitivity of this computation. The result therefore does
\emph{not} prove decoupling; it says that no coupling was detected in this small
specified family.

\section{Status of the claims}
\label{sec:status}

To keep the register unambiguous:
\begin{itemize}[leftmargin=1.6em]
\item \emph{Proved}: Lemmas~\ref{lem:counters}, \ref{lem:parity}, \ref{lem:w2};
Propositions~\ref{prop:closed}, \ref{prop:block}, \ref{prop:scale},
\ref{prop:info}, \ref{prop:offset}; Corollary~\ref{cor:bounds};
Theorems~\ref{thm:RH}, \ref{thm:gaps}.
\item \emph{Measured}, over $10^{4}\le n\le10^{6}$ with effective sample
$\approx1.7\times10^{4}$ where relevant: the law
$|E_{\mathrm{mid}}|\approx0.155\,|\Mert|^{3/2}\log n$; the failure of the
independent-gap model by a factor $0.853$ and its repair by the measured gap
covariance; approximate normality of the standardised residual; the
$0.9939$ packet-envelope correlation; the exact finite $\kappa$--$\rho$ transform (regenerated to double-precision
roundoff) and the reported finite-dimensional recovery of the first lags of $\rho$
from $\kappa$; the short-interval variance ratios; and the
null cross-coupling tests of Section~\ref{sec:coupling}.
\item \emph{Classical interpretation, not novelty}: the frequencies $\gamma_k$
in the $u=\log n$ spectrum of $M(e^u)e^{-u/2}$ are the modes supplied by the
explicit formula. The displayed peak matching is diagnostic only and has
frequency resolution about $0.91$.
\item \emph{Conjectural}: that $Z_{3/2}$ converges in distribution; that the
limit is Gaussian; that a central limit theorem holds for the antisymmetric
weighted gap sum of Theorem~\ref{thm:gaps}; or that the observed numerical
absence of cross-coupling persists asymptotically. None follows from the present
numerics.
\end{itemize}

\section{What the construction does and does not deliver}
\label{sec:deliverables}

It delivers an exact, elementary dictionary between $SS$ and $(\Mert,\Sq)$, a
two-sided transfer of the Mertens criterion (Theorem~\ref{thm:RH}), and a
residual whose size obeys a clean measured scaling law.

It does not deliver an independent arithmetic constraint on $\Mert(n)$, by
Proposition~\ref{prop:info}. For a transformed representation to become useful
it would have to produce at least one of:
\begin{enumerate}[label=(\roman*)]
\item a cancellation unavailable in the original variables;
\item a positivity or monotonicity principle strong enough to yield new bounds;
\item a spectral representation turning a Mertens bound into an operator inequality;
\item an inequality between independently controlled components;
\item a mechanism by which the prime side constrains the M\"obius side beyond
what $(\Mert,\Sq)$ already encodes.
\end{enumerate}
None of these has been extracted here. The measured law of
Section~\ref{sec:numeric} is of type (iv) in form only: it bounds
$E_{\mathrm{mid}}$ in terms of $|\Mert|$, which is the wrong direction to be
useful. The cross-coupling experiment tests whether the prime-gap fluctuation
might nevertheless provide an independently controlled handle on the Mertens
side; no such handle is detected, but the numerical sensitivity is only at the
$10^{-2}$ correlation scale and therefore does not close every conceivable
nonlinear or asymptotic coupling.

\section{A variant that is not an endpoint statistic (speculative)}
\label{sec:shared-counter}

The obstruction identified in Proposition~\ref{prop:info} is caused by the use
of two \emph{separate} ordinal counters, which destroys the ordering of the
$\pm1$ occurrences. The minimal modification that removes it is a
\emph{shared} counter:
\begin{equation}
\label{eq:shared-counter-definition}
\widetilde S(k):=\mob(k)\,p_{\Sq(k)},
\qquad
\widetilde{SS}(n):=\sum_{k\le n}\mob(k)\,p_{\Sq(k)} .
\end{equation}
Here the prime index depends on the position $k$ through the shared squarefree
counter rather than on a sign-specific ordinal. As a functional of a ternary history,
it is not an endpoint statistic: for example the histories $(+1,-1,+1)$ and
$(+1,+1,-1)$ have the same terminal counts $(M,Q)=(1,3)$ but shared-counter sums
$2-3+5=4$ and $2+3-5=0$, respectively. For the actual M\"obius sequence, finite
summation by parts with $w(k)=p_{\Sq(k)}$ gives exactly
\begin{equation}
\label{eq:abel-shared-counter}
\widetilde{SS}(n)=\Mert(n)w(n)-\sum_{k<n}\Mert(k)\bigl(w(k+1)-w(k)\bigr),
\end{equation}
The increments $w(k+1)-w(k)$ vanish when $\mob(k+1)=0$ and otherwise equal the
next prime gap, so the right-hand side samples the preceding values $M(k)$ along the
whole path. Moreover $Q(k)\sim6k/\pi^2$ and the prime number theorem imply
$w(k)\sim(6/\pi^2)k\log k$. This is offered as a non-endpoint direction, not as
a proved source of new bounds; asymptotically its bulk is close to a classical weighted
Mertens sum with weight $k\log k$.

\section{Novelty status}
\label{sec:novelty}

The three zeta ingredients are classical. The spectral occurrence of the
$\gamma_k$ in $M(e^u)e^{-u/2}$ is also classical explicit-formula content, and
$\kappa(h)$ is exactly a linear functional of the prime-gap autocorrelation
sequence, not a new GUE-type statistic. GUE concerns zero spacings; the prime
gaps entering Theorem~\ref{thm:gaps} belong to a different arithmetic
statistical problem.

A search did not locate the sequence $S$
($2,-2,-3,1,-5,3,-7,1,1,5,-11,1,\dots$) as a named sequence or standard
transform. Absence of a search result is weak evidence and the search was not
exhaustive; this is recorded as an observation about the formulation, not a
claim about its consequences, which by Proposition~\ref{prop:info} follow from
$(\Mert,\Sq)$. The genuinely non-endpoint variant in the preceding section remains a
direction rather than a result.

\section{Summary}
\label{sec:summary}

\begin{equation}
\label{eq:summary-map}
\mob \longrightarrow (N_-,N_0,N_+) \longrightarrow S \longrightarrow SS,
\qquad
N_\pm=\frac{\Sq\pm\Mert}2,\quad N_0=n-\Sq,\quad \Sq\equiv\Mert\ (2),
\end{equation}
\begin{equation}
\label{eq:summary-closed}
\boxed{\;
SS(n)=\PS\!\left(\frac{\Sq+\Mert}2\right)-\PS\!\left(\frac{\Sq-\Mert}2\right)+n-\Sq \;}
\end{equation}
\begin{equation}
\label{eq:summary-block}
\boxed{\;T(n)=SS(n)-n+\Sq(n)=\operatorname{sgn}(\Mert)\!\!\sum_{j\in W(n)}\!\!p_j,
\quad |W(n)|=|\Mert(n)|,\ \text{centre } \tfrac{\Sq(n)+1}2 \;}
\end{equation}
\begin{equation}
\label{eq:summary-scale}
\boxed{\;T(n)=\Mert(n)\,p_{\lfloor \Sq(n)/2\rfloor}\bigl(1+\theta(n)\bigr),\quad n\ge2,\quad \theta(n)\to0\;}
\end{equation}
\begin{equation}
\label{eq:summary-rh}
\boxed{\;\mathrm{RH}\iff T(n)=O_\varepsilon\bigl(n^{3/2+\varepsilon}\bigr)\ \ \forall\varepsilon>0\;}
\end{equation}
\begin{equation}
\label{eq:summary-residual}
\boxed{\begin{aligned}
T(n)&=\Mert(n)\,p^{*}(n)+E_{\mathrm{mid}}(n),\\
E_{\mathrm{mid}}(n)&=\operatorname{sgn}(\Mert)\sum_i w_i g_i,\\
\sum_i w_i^2&=\frac{h^3}{12}\bigl(1+O(h^{-1})\bigr).
\end{aligned}}
\end{equation}
\begin{equation}
\label{eq:summary-sd}
\boxed{\;\operatorname{SD}\bigl(E_{\mathrm{mid}}\bigr)\;=\;\sigma_g(n)\sqrt{w^{\!\top}\!R\,w}\,,
\qquad
|E_{\mathrm{mid}}(n)|\approx0.155\,|\Mert(n)|^{3/2}\log n\;}
\end{equation}
The gap representation and its prime-counting integral form are exact; the
$0.155$ scaling and distributional statements are measured over
$10^{4}\le n\le10^{6}$ and are not proved. The packets are the excursions of
$|M|$ amplified by the $h^{3/2}$ weight norm. Their zeta frequencies on
$u=\log n$ are classical explicit-formula modes, and the tested standardised
gap factor shows no detectable linear coupling to the sign or first zeta phase
of $M$ at current sensitivity.

The pointwise transform quantities $SS(n)$, $T(n)$, $E(n)$ and
$E_{\mathrm{mid}}(n)$ are deterministic functions of $(M(n),Q(n))$. This proves
that they contain no additional endpoint data. It does not, by itself, rule out an
analytic inequality obtained independently from the prime-gap representation; such an
inequality would have to supply leverage not already equivalent to a bound on $M$. The
finite coupling tests above find no such mechanism in the tested family. The shared-counter
variant of Section~\ref{sec:shared-counter} is genuinely non-endpoint as a functional
of sign histories and is therefore outside this endpoint obstruction.

\section*{References}
\label{sec:references}
\addcontentsline{toc}{section}{References}
\begin{thebibliography}{99}

\bibitem{Apostol1976}
T. M. Apostol,
\emph{Introduction to Analytic Number Theory},
Springer-Verlag, New York, 1976.

\bibitem{Cramer1936}
H. Cram\'er,
On the order of magnitude of the difference between consecutive prime numbers,
\emph{Acta Arith.} \textbf{2} (1936), 23--46.
DOI: 10.4064/aa-2-1-23-46.

\bibitem{MontgomerySoundararajan2004}
H. L. Montgomery and K. Soundararajan,
Primes in short intervals,
\emph{Comm. Math. Phys.} \textbf{252} (2004), 589--617.
DOI: 10.1007/s00220-004-1222-4; arXiv:math/0409258.

\bibitem{OdlyzkoTeRiele1985}
A. M. Odlyzko and H. J. J. te Riele,
Disproof of the Mertens conjecture,
\emph{J. Reine Angew. Math.} \textbf{357} (1985), 138--160.
DOI: 10.1515/crll.1985.357.138.

\bibitem{RieselGohl1970}
H. Riesel and G. G\"ohl,
Some calculations related to Riemann's prime number formula,
\emph{Math. Comp.} \textbf{24} (1970), 969--983.
DOI: 10.1090/S0025-5718-1970-0277489-3.

\bibitem{Titchmarsh1986}
E. C. Titchmarsh,
\emph{The Theory of the Riemann Zeta-Function},
2nd ed., revised by D. R. Heath-Brown, Clarendon Press, Oxford, 1986.

\end{thebibliography}

\end{document}
